% Nouveau EGA V (cote 25): archive cote 25, site page 80 (PDF page 81).
% typescript. Draft compared: none.
% AI-transcribed by Claude Opus 5.5 (claude-opus-5-5, adaptive thinking, effort high), Message Batches API, batch ega25b (msgbatch_01A7n5YuXubUPNK1zsepi8RG), 2026-11-13; not yet checked by eye.
\scanpage{25/80}
\mspage{\hand{75}}

\var{\margin{|}}{\placedmargin{|}}{marginal note placed}{ega25_25_080-v0001}m, en notant qu'à changement de corps de base près, $\underline{X}^{(m)}_{\hand{\unsure{\eta|\xi}}}$ s'obtient en
prenant une L' de codimension m-1 passant par x, une H de codi\var{\xout{me}}{}{typewriter strike-out}{ega25_25_080-v0002}nsion
1 passant par x, L' et H étant génériques indépendants pour ces proprié\var{\xout{x}}{}{typewriter strike-out}{ega25_25_080-v0003}
tés (i.e. en termes orthodoxes, on se place au point générique du \var{\xout{xx}}{}{typewriter strike-out}{ega25_25_080-v0004}
schéma des couples (L',H) \var{\xout{xxxxxx}}{}{typewriter strike-out}{ega25_25_080-v0005} ) et en prenant la section linéaire
de X par $L'\cap H$ : on peut commencer par prendre \var{\xout{\unsure{l}}}{}{typewriter strike-out}{ega25_25_080-v0006}a section par H, qui
est lisse par ce qui précède, puis par L', qui est lisse par l'hypothèse
de récurrence. Ce genre de raisonnement, déjà utilisé\var{\xout{e}}{}{typewriter strike-out}{ega25_25_080-v0007} pour généraliser
2.6. p.ex. aux sections linéaires de codimension m queconque, mérite
d'être explicité une bonne fois en général, pour pouvoir y référer sans
entrer \var{\xout{dans}}{}{typewriter strike-out}{ega25_25_080-v0008} chaque fois dans les détails un peu lourds d'une présentati-%
on en forme.

Reste à prouver le corollaire annoncé de 8.8. dans le cas m=1 .
Lorsque \var{\xout{dim}}{}{typewriter strike-out}{ega25_25_080-v0009} codim T $\geqslant$ 2, comme d'autre part l'hyperplan \var{\inserted{Q}}{Q}{insertion}{ega25_25_080-v0010} de $\overset{\hand{\vee}}{P}$ des
$H_{\hand{\xi}}$ tels que $x\in H_{\hand{\xi}}$ est de codimension 1, son point générique \var{\inserted{$\eta$}}{$\eta$}{insertion}{ega25_25_080-v0011} ne peut
être élement de T, et on gagne. Dans le cas codim T = 1, comme T est
irréductible, on ne peut avoir $\hand{\eta}\in T$ que si Q=T ,\var{\xout{xxxxxxxxxxxxx}}{}{typewriter strike-out}{ega25_25_080-v0012}
i.e. en termes géométriques \var{\xout{xxxx}}{}{typewriter strike-out}{ega25_25_080-v0013} \var{\inserted[typed]{(supposant k alg clos, ce qui est loisible)}}{(supposant k alg clos, ce qui est loisible)}{insertion}{ega25_25_080-v0014} pour tout $z\in X(\unsure{k})$\var{\xout{xxx}}{}{typewriter strike-out}{ega25_25_080-v0015}l'espace tangent
de X en z (ou plutôt son image par $f'_{z}$ ) passe\var{\xout{x}}{}{typewriter strike-out}{ega25_25_080-v0016} par x. Montrons que
ceci ne peut se présenter lorsque $\underline{X}^{\mathrm{sing}}\hand{\longrightarrow}T$ est gén étale, i.e.
lorsqu'on est sous les conditions de 8.8., sauf si f(X) = \hand{\{}x\hand{\}} donc X
de dimension 0. En effet, 8.7. c) (qui exprime essentiellement la
symétrie \var{\xout{xxxxx}}{}{typewriter strike-out}{ega25_25_080-v0017} dans la relation entre X et sa \hand{``}duale\hand{''} \var{\inserted{T}}{T}{insertion}{ega25_25_080-v0018}\hand{)} implique alors que
pour presque tout point $z\in X(k)$, f(z) est orthogonal à l'espace tangent
à T en un certain point,\var{\xout{xx}}{}{typewriter strike-out}{ega25_25_080-v0019} donc (puisque T=Q) \var{\inserted{orthogonal}}{orthogonal}{insertion}{ega25_25_080-v0020} à Q, d'où f(z)=x, d'où
f(X) =\hand{\{}x\hand{\}} . Cela achève la démonstration de nos remords, donc de 13.2.

\var{\margin{80}}{\placedmargin{80}}{marginal note placed}{ega25_25_080-v0021}
